\section{算力与成本：从训练可行到产品可用}

前面已经说明训练阶段会改变主要稀缺资源，本节把这种差异写成可核算的成本模型。这里不追求给所有项目套一个美元单价，而是区分一次性的训练投入、持续发生的推理费用、评估环境成本与运行维护成本，并解释为什么 Agent 的长链路会让后两项快速上升。

\subsection{课程中的成本讨论框架}
本小节承接官方 slide 22 的数量级估算，把课堂中的 FLOPs、GPU-hours 和推理链路转换成项目台账。阅读时应把每个数字当作依赖硬件、利用率和任务分布的变量，而不是固定常数；真正要建立的是“假设可见、公式可复算、版本可比较”的成本纪律。
讲者在多个位置给出 ``compute cost around ...'' 的数量级提示，核心目的是让研究决策与资源现实对齐。即便某条路径在离线实验有效，也必须评估其训练时间、并行开销、推理成本和部署成本。

\begin{importantbox}{训练成本与推理成本必须分开记账}
课程明确指出，很多实验只优化训练资源，却忽略 inference cost。对于 Agent 产品，推理侧往往是持续性支出，且会随着任务链路拉长而成倍增长，因此不能只看训练阶段的 ``一次性'' 成本。
\end{importantbox}

\subsection{成本分解模型}
可把总成本粗略拆为：
\[
C_{\text{total}} = C_{\text{data}} + C_{\text{train}} + C_{\text{eval}} + C_{\text{inference}} + C_{\text{ops}}
\]
其中最容易被低估的是 $C_{\text{eval}}$ 与 $C_{\text{inference}}$。Agent 场景下，评估本身可能需要复杂环境复现，推理则包含多轮调用、工具耗时和失败重试。

\begin{knowledgebox}{对齐研发与产品的成本仪表盘}
建议至少监控以下指标：
\begin{itemize}
    \item 每次有效策略改进的平均 token 成本；
    \item 每个成功任务的平均 rollout 长度；
    \item 失败重试率与重试引入的额外推理费用；
    \item 评估覆盖率与评估单次执行成本。
\end{itemize}
\end{knowledgebox}

\subsection{本章小结}
成本不是财务层面的后置问题，而是训练策略选择的一阶约束。忽略推理和评估成本，会让 ``研究上有效'' 的方案无法产品化。


